\section*{Chapitre \ref{ch:m1:shares-corporate-actions-indices} --- Actions, opérations sur titres et indices}

\begin{solution}{exo:m1:shares-corporate-actions-indices:1}
52~milliards de dollars~; ajustée du \omterm{def:m1:shares-corporate-actions-indices:float}{flottant}, 36.4~milliards de dollars.
\end{solution}

\begin{solution}{exo:m1:shares-corporate-actions-indices:2}
Chaque action devient 1.5 actions~: cours d'ouverture $240/1.5 = \$160$, facteur
$2/3$, position de 750 actions (valant toujours 120\,000~\$).
\end{solution}

\begin{solution}{exo:m1:shares-corporate-actions-indices:3}
Cours ex 78.80~\$~; $f = 78.80/80.00 = 0.985$~; $76.00 \times 0.985 = \$74.86$.
\end{solution}

\begin{solution}{exo:m1:shares-corporate-actions-indices:4}
$\mathrm{TERP} = (7\times25 + 2\times16)/9 = 23.00$~; un droit vaut
$23 - 16 = 7.00$~; $f = 0.92$. Un détenteur de 700 actions reçoit des droits sur 200
nouvelles actions, les vend pour 1\,400~\euro, et détient $700 \times 23 = 16\,100$~:
17\,500~\euro, comme auparavant.
\end{solution}

\begin{solution}{exo:m1:shares-corporate-actions-indices:5}
Niveau $500/0.5 = 1\,000$. Après la \omterm{def:m1:shares-corporate-actions-indices:split}{division d'actions}, le troisième cours est 110 et la somme
280, donc $\Delta' = 280/1\,000 = 0.28$. Les pondérations passent de 10~\%, 24~\%, 66~\% à
17.9~\%, 42.9~\%, 39.3~\%~: l'indice est désormais un autre portefeuille, par une décision
du conseil d'administration de la société et non du fournisseur d'indices.
\end{solution}

\begin{solution}{exo:m1:shares-corporate-actions-indices:6}
Facteur standard~: $48.45/47.50 - 1 = +2.00\%$. Rendement total~: $(48.45 +
2.50)/50 - 1 = +1.90\%$. Écart de 0.10~\%, égal à $\varepsilon D/P = 0.02
\times 0.05$.
\end{solution}

\begin{solution}{exo:m1:shares-corporate-actions-indices:7}
(a) \omterm{def:m1:shares-corporate-actions-indices:share}{Capitalisation boursière} 5~milliards de dollars, diviseur 5\,000\,000. (b) La \omterm{def:m1:shares-corporate-actions-indices:share}{capitalisation boursière} tombe à
$2.0 + 2.4 = 4.4$ milliards de dollars aux cours inchangés~: diviseur 4\,400\,000, niveau
toujours 1\,000. (c) $(100\times22 + 40\times57)/4.4 = 1\,018.2$.
\end{solution}

\begin{solution}{exo:m1:shares-corporate-actions-indices:8}
D'abord, le cours ajusté d'une action il y a vingt ans est son cours brut multiplié par
chacun des facteurs postérieurs~: une action qui a depuis subi de nombreuses divisions affiche un ancien
cours ajusté inférieur à 5~\$ alors qu'elle se négociait à 80~\$ ce jour-là. La règle
retient donc, dans le passé, précisément les actions qui sont ensuite montées assez
pour être divisées~: le futur est dans le filtre. Deuxièmement, une liste téléchargée du jour
des 3\,000 actions cotées ne contient que des survivants~; les valeurs à bas cours tombées à
zéro sont absentes. Les deux erreurs avantagent la règle~; le test exige les cours bruts
à chaque date et un univers à chaque date.
\end{solution}

\begin{solution}{pb:m1:shares-corporate-actions-indices:1}
\textbf{1.} P 20.0, Q 24.0, R 6.0, S 27.0~: 77.0~milliards de dollars.
\textbf{2.} $77\times10^9/2\,500 = 30\,800\,000$.
\textbf{3.} 26.0~\%, 31.2~\%, 7.8~\%, 35.1~\%.
\textbf{4.} $0.01 \times 24\times10^9/30.8\times10^6 = 7.8$ points.
\textbf{5.} Cours 50.00, 600 millions d'actions~; la capitalisation et donc
le diviseur sont inchangés.
\textbf{6.} La capitalisation baisse de $500 \times 1.00 = 0.5$ milliard de dollars~: l'indice
est $76.5\times10^9/30.8\times10^6 = 2\,483.77$, en baisse de 16.23 points d'indice. Aucun
ajustement~: dans un indice de prix, un \omterm{def:m1:shares-corporate-actions-indices:dividend}{dividende} ordinaire est un mouvement de cours.
\textbf{7.} Il réintègre les 16.23 points d'indice~: 2\,500.00.
\textbf{8.} L'\omterm{def:m1:shares-corporate-actions-indices:tri}{indice de rendement total} (net de la retenue à la source si le fonds
la subit)~; face à l'indice de prix, le fonds paraîtrait surperformer de
chaque \omterm{def:m1:shares-corporate-actions-indices:dividend}{dividende}.
\textbf{9.} $\mathrm{TERP} = (4\times12 + 8)/5 = 11.20$~; 1\,250 millions
d'actions.
\textbf{10.} $11.20 \times 1\,250 \times 0.5 = 7.0$ milliards de dollars~; la \omterm{def:m1:shares-corporate-actions-indices:share}{capitalisation
boursière} $76.5 - 6.0 + 7.0 = 77.5$ milliards de dollars.
\textbf{11.} $\Delta' = 30.8 \times 77.5/76.5 = 31\,202\,614$~; le niveau
reste 2\,483.77.
\textbf{12.} La pondération de R augmente parce que la société est plus grande~: le fonds
souscrit à ses droits (ou vend une partie de ses droits et achète des actions pour le même
effet), en investissant des liquidités nouvelles dans R, financées par la vente d'un peu de tout
le reste.
\textbf{13.} $77.5 - 27.0 + 13.5 = 64.0$ milliards de dollars~; $\Delta_{\text{nouveau}} = 31\,202\,614
\times 64/77.5 = 25\,767\,320$.
\textbf{14.} P 30.5~\%, Q 37.5~\%, R 10.9~\%, T 21.1~\%.
\textbf{15.} S représente $27/77.5 = 34.8\%$ du fonds~: il reçoit 697~millions de dollars
en numéraire de la part de l'acquéreur et doit l'investir dans le nouvel indice, dont
422~millions de dollars dans T~; le reste, soit 275~millions de dollars, est réparti entre P, Q et R.
\textbf{16.} $422/150 = 2.8$ jours de volume~; coût $0.7\times0.02\times
\sqrt{2.8} = 2.3\%$, soit environ 9.9~millions de dollars --- pour un seul fonds.
\textbf{17.} Les traders qui ont acheté T lors de l'annonce du changement, et
les arbitragistes qui achètent pendant la semaine afin de revendre au
cours de clôture effectif~: le coût du fonds est leur revenu (livre~8).
\textbf{18.} Les trackers ont besoin de temps pour organiser la transaction, et un
changement non annoncé ferait de la propre fuite du fournisseur la seule information~: le délai
échange un coût connu (front-running) contre l'équité et une exécution ordonnée.
Certains fournisseurs prévoient un préavis plus long et une mise en œuvre échelonnée pour répartir
l'impact.
\textbf{19.} \textbf{25\,767\,320.}
\textbf{20.} Modifier le niveau de l'indice.
\end{solution}

\begin{solution}{iq:m1:shares-corporate-actions-indices:1}
Une \omterm{def:m1:shares-corporate-actions-indices:split}{division d'actions} 2 pour 1 (ou un gros \omterm{def:m1:shares-corporate-actions-indices:dividend}{dividende} exceptionnel ou une \omterm{def:m1:shares-corporate-actions-indices:spinoff}{scission}). Le pipeline doit
ingérer les opérations sur titres provenant d'une source de référence, conserver les cours bruts et
les opérations séparément, fournir des séries rétro-ajustées pour le calcul des rendements, et
ajuster les positions ouvertes, les ordres ouverts et les cours de référence au même moment.
Un mouvement important de nuit sans opération enregistrée doit déclencher une alerte avant
qu'elle n'atteigne une stratégie.
\iqlookfor{les positions et les ordres, et pas seulement l'historique des cours.}
\end{solution}

\begin{solution}{iq:m1:shares-corporate-actions-indices:2}
Du fait du \omterm{def:m1:shares-corporate-actions-indices:dividend}{dividende}, en l'absence d'impôts et de nouvelles, par l'absence d'arbitrage entre le cours
cum de clôture et le cours ex d'ouverture. Moins quand les \omterm{def:m1:shares-corporate-actions-indices:dividend}{dividendes} sont plus imposés que les plus-values en capital
pour l'investisseur marginal, et les baisses mesurées sont aussi brouillées par le
mouvement de marché de nuit et par le pas de cotation.
\iqlookfor{l'argument d'arbitrage, puis l'effet de clientèle fiscale.}
\end{solution}

\begin{solution}{iq:m1:shares-corporate-actions-indices:3}
Brut pour tout ce que voit la bourse~: cours limites, règles de tick, simulation
d'exécution, frais. Ajusté pour les rendements et les indicateurs à travers les \omterm{def:m1:shares-corporate-actions-indices:dividend}{dates de détachement}. Brut
utilisé pour les rendements~: une \omterm{def:m1:shares-corporate-actions-indices:split}{division d'actions} apparaît comme une journée à $-90\%$ et le modèle de risque
explose. Ajusté utilisé pour les ordres ou les filtres~: un backtest «~achète~» à des cours
qui n'ont jamais été échangés, ou sélectionne des actions en utilisant une information issue de \omterm{def:m1:shares-corporate-actions-indices:split}{divisions d'actions} postérieures.
\iqlookfor{le biais de regard vers l'avenir dans les cours ajustés.}
\end{solution}

\begin{solution}{iq:m1:shares-corporate-actions-indices:4}
Une pondération par la capitalisation vaut cours multiplié par nombre d'actions, quantité qu'une \omterm{def:m1:shares-corporate-actions-indices:split}{division d'actions} laisse
inchangée. Une pondération par le cours ne tient que du cours seul, que la division divise~; le
\omterm{def:m1:shares-corporate-actions-indices:divisor}{diviseur d'indice} rétablit le niveau mais ne peut pas rétablir les pondérations.
\iqlookfor{une phrase chacune~; bonus pour «~le diviseur fixe le niveau, pas
la composition~»}
\end{solution}

\begin{solution}{iq:m1:shares-corporate-actions-indices:5}
Stocker~: cours et volumes bruts tels qu'échangés, immuables~; opérations sur titres avec \omterm{def:m1:shares-corporate-actions-indices:dividend}{date de
détachement}, conditions et horodatage de l'annonce~; historique des identifiants. Calculer à la
lecture~: \omterm{def:m1:shares-corporate-actions-indices:factor}{facteurs d'ajustement} et séries ajustées à une date de référence donnée, de sorte que
les requêtes à date donnée soient possibles et qu'une correction tardive d'une seule opération corrige
toutes les séries d'un coup. Ne jamais stocker les cours ajustés comme données primaires~: ils deviennent
obsolètes à chaque nouvelle opération et masquent la convention utilisée.
\iqlookfor{requêtes à date de référence et immuabilité.}
\end{solution}

\begin{solution}{iq:m1:shares-corporate-actions-indices:6}
Non. L'effet mécanique d'une émission 1 pour 3 à une décote de 40~\% est un TERP
inférieur de 10~\% au cours cum, et il apparaît à la \emph{\omterm{def:m1:shares-corporate-actions-indices:dividend}{date de détachement}}, le détenteur
étant compensé par les droits. Une baisse à l'\emph{annonce} relève de l'information~: l'entreprise a besoin de
capitaux, peut-être en urgence, et la direction a choisi d'émettre des actions à ce prix~; les détenteurs qui ne peuvent pas souscrire doivent vendre
leurs droits sur un marché qui le sait. Une baisse de 12~\% à l'annonce
est une perte de valeur~; la baisse ultérieure de 10~\% à la \omterm{def:m1:shares-corporate-actions-indices:dividend}{date de détachement}
n'en est pas une.
\iqlookfor{séparer l'effet de l'annonce de l'ajustement
ex-droits.}
\end{solution}
